\documentclass[letter, 10pt]{article}
\usepackage[utf8]{inputenc}
\usepackage[spanish]{babel}
\usepackage{amsfonts}
\usepackage{amsmath}
\usepackage[dvips]{graphicx}
\usepackage{url}
\usepackage[top=3cm,bottom=3cm,left=3.5cm,right=3.5cm,footskip=1.5cm,headheight=1.5cm,headsep=.5cm,textheight=3cm]{geometry}

\begin{document}
\title{Inteligencia Artificial \\ \begin{Large}Informe Final: Problema ARSP\end{Large}}
\author{Sergio Rojas Gutiérrez}
\date{\today}
\maketitle


%--------------------No borrar esta secci\'on--------------------------------%
\section*{Evaluaci\'on}

\begin{tabular}{ll}
Código Fuente (10\%): &  \underline{\hspace{2cm}}\\
Resumen (2\%):  & \underline{\hspace{2cm}} \\
Introducci\'on (3\%):  & \underline{\hspace{2cm}} \\
Representaci\'on (10\%):  & \underline{\hspace{2cm}} \\
Descripci\'on del algoritmo (20\%):  & \underline{\hspace{2cm}} \\
Experimentos (10\%):  & \underline{\hspace{2cm}} \\
Resultados (30\%):  & \underline{\hspace{2cm}} \\
Conclusiones (10\%): &  \underline{\hspace{2cm}}\\
Bibliograf\'ia (5\%): & \underline{\hspace{2cm}}\\
 &  \\
\textbf{Nota Final (100)}:   & \underline{\hspace{2cm}}
\end{tabular}
%---------------------------------------------------------------------------%

\begin{abstract}
Este informe aborda el Aircraft Runway Scheduling Problem (ARSP), un desafío logístico crítico para descongestionar el tráfico aéreo. El objetivo es optimizar secuencias y tiempos de arribo para minimizar costos por demoras, respetando ventanas de vuelo y separaciones de seguridad. Se propone e implementa un algoritmo metaheurístico de búsqueda local (Hill Climbing con Best Improvement y Restarts). Evaluado en pistas únicas y múltiples, el método demuestra empíricamente una superioridad absoluta frente al orden de llegada tradicional (FCFS), mitigando penalizaciones velozmente.
\end{abstract}

\section{Introducción}

La gestión eficiente del espacio aéreo terminal representa uno de los desafíos logísticos más críticos en la aviación moderna. El constante incremento en el volumen de tráfico aéreo global exige sistemas de control automatizados capaces de tomar decisiones óptimas en fracciones de segundo, garantizando la seguridad operacional sin comprometer la fluidez de las plataformas aeroportuarias. 

El propósito fundamental de este documento es presentar el diseño, implementación y evaluación experimental de un enfoque metaheurístico basado en trayectorias para resolver el Aircraft\textit{} Runway Scheduling Problem (ARSP). Específicamente, se analiza el comportamiento del algoritmo Hill Climbing con Mejor Mejora y Reinicios (HC-BI) bajo configuraciones estáticas de pista única y múltiples pistas de aterrizaje.

El problema analizado consiste en determinar una secuencia de arribo y asignar horarios de aterrizaje precisos para un conjunto de aeronaves en cola. Esta asignación busca minimizar los costos asociados a las penalizaciones por desvíos temporales (adelantos o retrasos) respecto al tiempo objetivo preferido de cada operador, sujeto a dos restricciones duras e inquebrantables: las ventanas de tiempo operativas de cada avión y las separaciones mínimas impuestas por la estela de turbulencia asimétrica generada entre aeronaves consecutivas. 

La principal motivación detrás de este estudio radica en el severo impacto económico y ambiental que conllevan las ineficiencias en la secuenciación de vuelos, tales como el gasto innecesario de combustible y las demoras operativas en cadena. Dado que el ALP pertenece a la clase de problemas NP-Hard, los métodos de optimización matemática exacta colapsan ante la explosión combinatoria en instancias de escala real. Esto justifica el desarrollo de propuestas de solución heurísticas como la presentada en este informe, la cual utiliza una representación eficiente basada en arreglos de permutación explorados de forma exhaustiva mediante un operador de movimiento de intercambio (\textit{Swap}). La estrategia combina la explotación intensiva del entorno local con mecanismos de diversificación estocástica mediante reinicios paramétricos y asignaciones voraces (\textit{Greedy}) para el escenario de pistas múltiples.

Para ofrecer una exposición clara y ordenada del trabajo realizado, el resto de este informe se estructura de la siguiente manera:
\begin{itemize}
    \item \textbf{Sección 2 (Definición del Problema):} Desglosa los componentes operativos del ARSP, sus variables de decisión y las variantes metodológicas documentadas.
    \item \textbf{Sección 3 (Estado del Arte):} Revisa la evolución histórica de la literatura, desde los primeros modelos MIP clásicos hasta las tendencias de flujo lineal y aprendizaje profundo.
    \item \textbf{Sección 4 (Modelo Matemático):} Detalla las formulaciones de programación entera mixta de Beasley y el enfoque de flujo de costo mínimo para asignación de tiempos.
    \item \textbf{Sección 5 (Representación y Restricciones):} Explica la estructura de datos del algoritmo y los mecanismos implementados para garantizar la factibilidad de las soluciones.
    \item \textbf{Sección 6 (Descripción del Algoritmo):} Expone la lógica procedimental y el pseudocódigo del Hill Climbing implementado.
    \item \textbf{Sección 7 (Experimentos):} Describe el entorno de hardware, las instancias de prueba seleccionadas (Airland y FPT) y los criterios de variación paramétrica.
    \item \textbf{Sección 8 (Resultados y Análisis):} Contrasta empíricamente el rendimiento del algoritmo frente a la línea base FCFS, analizando su escalabilidad en múltiples pistas y el \textit{trade-off} computacional de los hiperparámetros.
    \item \textbf{Sección 9 (Conclusiones):} Sintetiza las principales deducciones e implicancias ingenieriles del estudio.
    \item \textbf{Sección 10 (Uso de LLMs):} Declara con transparencia el uso y validación de herramientas de asistencia basada en IA.
\end{itemize}

\section{Definici\'on del Problema}

El \textbf{Problema de Aterrizaje de Aeronaves} (ALP) se define como la tarea de programar el uso de la(s) pista(s) de aterrizaje en un aeropuerto para un conjunto de aviones que se aproximan a la zona de control.

Fundamentalmente, el problema busca decidir una secuencia de aterrizaje y asignar tiempos espec\'ificos para cada operaci\'on, esto con el objetivo de garantizar la seguridad de las aeronaves mientras se optimiza la eficiencia del flujo a\'ereo\cite{beasley2000}, lo que a su vez disminuye los gastos para el aeropuerto, el consumo de combustible y la huella de carbono asociada a mantener las aeronaves en el cielo.

\subsection{Componentes del problema}
En l\'ineas generales, el problema se estructura bajo los siguientes componentes operativos:
\begin{itemize}
    \item Objetivo: La meta \'ultima es minimizar el costo total por desviaci\'on respecto a los horarios ideales de llegada. Esto busca reducir tanto el aterrizaje prematuro (\textit{earliness}), que puede generar congesti\'on en las plataformas de desembarque, como el aterrizaje tard\'io (\textit{tardiness}), que genera retrasos en cadena, molestia en la clientela y costos operativos adicionales. \cite{beasley2000}
    \item Variables de Decisi\'on: El elemento esencial a determinar es el momento preciso de aterrizaje para cada avi\'on. A pesar de que la secuencia de aterrizaje (el orden de estos mismos) es un resultado de estos tiempos, la decisi\'on algor\'itmica implica encontrar un momento preciso dentro de un margen permitido para cada aeronave.
    \item Restricciones: Todo lo que se ha mencionado est\'a sujeto a dos tipos de limitaciones cr\'iticas.
    \begin{itemize}
        \item Ventanas de Tiempo: Estas ventanas definen un rango entre el momento m\'as temprano posible (dependiendo de la velocidad m\'axima del avi\'on) y el m\'as tard\'io posible (dependiendo del consumo de combustible) en el que se puede realizar el aterrizaje.
        \item Separaciones de Seguridad: Estas separaciones exigen un intervalo de tiempo m\'inimo entre aterrizajes consecutivos, estos intervalos dependen del tiempo necesario para que se disipe la estela de turbulencia generada por el aterrizaje m\'as reciente. Estos tiempos variar\'an seg\'un la categor\'ia de peso de cada aeronave involucrada.
    \end{itemize}
\end{itemize}

\subsection{Problemas Relacionados y Variantes}
El ALP no es un problema aislado, este forma parte de una familia de problemas de optimizaci\'on relacionados a la gesti\'on del tr\'afico a\'ereo (ATM - \textit{Air Traffic Management}). La estructura del mismo es an\'aloga al problema de secuenciamiento en una sola m\'aquina (\textit{Single Machine Scheduling}) con fechas de disponibilidad y penalizaciones por incumplimiento\cite{beasley2000}
. Asimismo, presenta una relaci\'on directa con el \textbf{Problema del Viajero Asim\'etrico (ATSP)}, donde los aviones representan los nodos, mientras que los tiempos de separaci\'on act\'uan como los costos de transici\'on dependientes del orden. \cite{prakash2018}


El problema cuenta con variantes que aumentan la complejidad, estas variantes nacen seg\'un necesidades del aeropuerto estudiado:
\begin{itemize}
    \item \textbf{ALP Est\'atico vs. Din\'amico:} Mientras que la variante est\'atica considera un conjunto de aviones fijos conocidos de antemano\cite{beasley2000}, la variante din\'amica gestiona la aparici\'on continua de nuevos aviones en tiempo real, esto requiere algoritmos de re-planificaci\'on constante o reactivos \cite{ikli2021}
    \item \textbf{Pista \'Unica vs. M\'ultiples Pistas:} El problema puede escalar a la gesti\'on simult\'anea de varias pistas, lo que a\~nade la variable de asignar cada avi\'on a una pista espec\'ifica para balancear la carga de trabajo\cite{beasley2000}.
    \item \textbf{Modo Mixto:} Algunas variantes consideran pistas que operan tanto para despegues como para aterrizajes de forma intercalada, lo que introduce restricciones adicionales de seguridad, esto por la ocupaci\'on de pista, o pistas, para diferentes tipos de maniobras\cite{bennell2011}.
\end{itemize}



\section{Estado del Arte}
El estudio del ALP ha evolucionado a pasos agigantados durante el \'ultimo medio siglo. Durante el transcurso de esta secci\'on se presenta una revisi\'on exhaustiva de la literatura, seccionada para responder a las interrogantes esenciales sobre la historia inicial y contempor\'anea de este desaf\'io log\'istico.

\subsection{Origen y Formalizaci\'on del Problema}
El inter\'es por la optimizaci\'on del flujo a\'ereo terminal tiene sus primeras ra\'ices documentadas en la d\'ecada de 1950. Se considera que el trabajo m\'as temprano en este \'ambito fue realizado por Blumstein (1959) \cite{blumstein1959}, quien describi\'o las primeras metodolog\'ias para la secuenciaci\'on y programaci\'on de aeronaves. 

Posteriormente, el salto hacia el modelamiento t\'ecnico moderno se consolid\'o con el trabajo pionero de Dear (1976) \cite{dear1976} en el MIT. En su tesis, Dear demostr\'o que la disciplina \textit{First-Come, First-Served} (FCFS) era ineficiente e introdujo el concepto de \textbf{\textit{Constrained Position Shifting} (CPS)} para garantizar equidad operativa. Sin embargo, Dear advirti\'o que optimizar secuencias en tiempo real era computacionalmente inviable para la tecnolog\'ia de su \'epoca.

Durante la d\'ecada de 1970, cercano al final de la misma, se estableci\'o la categoría de problema \textbf{NP}, con el documento m\'as conocido de Garey et. al \cite{garey1979}. Esta estandarizaci\'on fue un hito cr\'itico.

Finalmente, la estandarizaci\'on del ALP est\'atico moderno lleg\'o con Beasley et al. (2000) \cite{beasley2000}, quienes establecieron la formulaci\'on de programaci\'on entera mixta (MIP) y la librer\'ia de instancias de prueba (\textit{OR-Library}) que sigue siendo el est\'andar global de evaluaci\'on.


\subsection{M\'etodos Completos (Exactos)}
En los inicios de la investigaci\'on el foco estuvo en los m\'etodos completos, en busca de los \'optimos te\'oricos globales.

Las t\'ecnicas predominantes fueron la \textbf{Programaci\'on Din\'amica (DP)} y los m\'etodos de ramificaci\'on y acotamiento (\textit{Branch and Bound}) basados en modelos MIP \cite{beasley2000}

Para abordar la asimetr\'ia de las clases de aeronaves, investigadores como Prakash et al. (2018) \cite{prakash2018} introdujeron algoritmos de divisi\'on de datos (\textit{data-splitting}), relacionando el ALP con el Problema del Viajero Asim\'etrico con Ventanas de Tiempo (ATSP-TW). No obstante, la literatura llega a una conclusi\'on clara: los m\'etodos completos colapsan ante la explosi\'on combinatoria. \textit{Solvers} comerciales como CPLEX son incapaces de resolver instancias de m\'as de 50 aviones en tiempos operativos viables.

\subsection{Heur\'isticas y M\'etodos Populares}
Dada la naturaleza prohibitiva de los m\'etodos exactos para la operaci\'on en tiempo real, la literatura se ha centrado en el desarrollo de enfoques aproximados. 

La estrategia m\'as popular y exitosa consiste en descomponer el ALP en una cadena de sub-problemas m\'as sencillos, separando la decisi\'on del \textbf{orden de la secuencia} de la \textbf{asignaci\'on de los tiempos} de aterrizaje\cite{ikli2021}.

Hist\'oricamente, una vez determinado un orden tentativo de aviones, el m\'etodo predominante para calcular los tiempos de aterrizaje consist\'ia en resolver un sub-modelo de \textbf{Programaci\'on Lineal (LP)}. 

Aunque este enfoque garantiza la asignaci\'on \'optima de tiempos para esa secuencia espec\'ifica, requiere llamar a un \textit{solver} (como CPLEX o Gurobi) miles de veces durante la ejecuci\'on de una metaheur\'istica, lo que ralentiza dr\'asticamente el proceso. 

Alternativamente, se han utilizado algoritmos con complejidad cuadr\'atica $\mathcal{O}(n^2)$ que, si bien son m\'as ligeros que un modelo LP completo, siguen presentando limitaciones de escalabilidad cuando el n\'umero de aeronaves ($n$) crece significativamente.

Como respuesta a este cuello de botella computacional, Cao y Xu (2024)\cite{cao2024} introdujeron un algoritmo disruptivo que redefine la asignaci\'on de tiempos. 

Su enfoque consiste en transformar el problema de tiempos en una red de \textbf{flujo de costo m\'inimo} (\textit{min-cost flow}). El aporte fundamental de este m\'etodo es su eficiencia temporal: logra resolver la asignaci\'on \'optima en un tiempo casi lineal $\mathcal{O}(n \log n)$. 

Esta mejora permite que los algoritmos de b\'usqueda contempor\'aneos eval\'uen una cantidad de secuencias candidatas mucho mayor en el mismo intervalo de tiempo. En consecuencia, la combinaci\'on de una heur\'istica de b\'usqueda de secuencia (como una b\'usqueda local o algoritmos gen\'eticos) con el m\'etodo de flujo de Cao y Xu se ha consolidado como la tendencia actual para alcanzar soluciones de alta calidad en entornos operativos donde cada milisegundo es cr\'itico.

\subsection{Representaciones y Tipos de Movimientos}

Cuando empezamos a usar heur\'isticas computacionales y m\'etodos incompletos necesitamos una \textbf{representaci\'on} que est\'e a la altura para asegurar la eficacia.
La estructura de datos m\'as exitosa y estandarizada es el \textit{arreglo de permutaci\'on} (un vector donde cada \'indice representa el orden de llegada de un avi\'on espec\'ifico).

Para navegar y explorar el espacio de b\'usqueda desde una soluci\'on inicial, los algoritmos aplican diversos 
\textbf{tipos de movimientos} (\textit{move types}). 
Los que han demostrado mayor efectividad para el ALP incluyen:
\begin{itemize}
    \item \textbf{Intercambio (\textit{Swap}):} Intercambiar la posici\'on de dos aeronaves en la secuencia.
    \item \textbf{Inserci\'on (\textit{Shift}):} Tomar una aeronave y desplazarla a una nueva posici\'on, desplazando al resto de la cadena.
    \item \textbf{Inversi\'on (\textit{2-opt}):} Invertir el orden de una subsecuencia completa de aterrizajes.
\end{itemize}

\subsection{Los Mejores Algoritmos}
Con el avance del tiempo, la tecnolog\'ia y el conocimiento, las heur\'isticas simples dieron pie a las \textbf{Metaheur\'isticas}. Si vamos por la categor\'ia de la inteligencia bio-inspirada, el Algoritmo de Colonia de Hormigas (ACO) acoplado con modelado de v\'ortice demostr\'o superar a los \textit{solvers} matem\'aticos tradicionales \cite{xu2017}

Sin embargo, los mejores algoritmos creados hasta la fecha, y los que ostentan los r\'ecords mundiales en las instancias de Beasley, son los \textbf{algoritmos h\'ibridos}. Hammouri et al. (2020) \cite{hammouri2020} desarrollaron el algoritmo ISA (\textit{Iterated Simulated Annealing}), el cual mezcla la B\'usqueda Local Iterada (ILS) para diversificar la b\'usqueda usando los movimientos antes mencionados, con el Enfriamiento Simulado (SA) para escapar de \'optimos locales, logrando resultados sin precedentes.

\subsection{Tendencias Actuales}
La investigaci\'on contempor\'anea del ALP se dirige hacia dos grandes fronteras:
\begin{enumerate}
    \item \textbf{Optimizaci\'on Multi-objetivo:} Se ha pasado de solo minimizar retrasos a modelos que incluyen la reducci\'on de costos de combustible, el balance de las pistas y la equidad entre aerol\'ineas, requiriendo algoritmos que generen frentes de Pareto \cite{zhang2020}
    \item \textbf{Inteligencia Artificial y Aprendizaje Profundo:} La tendencia m\'as disruptiva es el uso de Redes Neuronales de Grafos (GNN) y Aprendizaje por Refuerzo Profundo (DRL). Maru (2025)\cite{maru2025} introdujo una nueva \textit{representaci\'on basada en grafos} donde los aviones son nodos y las separaciones son aristas temporales. Este algoritmo es capaz de reducir el tiempo de c\'omputo en un 99.95\% comparado con MIP, generando soluciones en menos de 1 segundo para entornos reactivos.
\end{enumerate}

Como conclusi\'on del an\'alisis a la literatura que rodea este problema, para un entorno est\'atico y resoluci\'on pr\'actica, el dise\~no de 
algoritmos de b\'usqueda heur\'istica implementados en lenguajes de alto rendimiento (como C o C++) se mantiene 
como el enfoque metodol\'ogico m\'as equilibrado entre factibilidad de desarrollo y excelencia en 
tiempos operacionales.

\section{Modelo Matem\'atico}

En esta secci\'on se presentan las dos formulaciones fundamentales que sustentan este estudio. Primero, el modelo cl\'asico unificado de Beasley, que abarca la asignaci\'on de pistas y secuencias. Segundo, el avance de 2024 de Cao y Xu, que optimiza exclusivamente la asignaci\'on de tiempos para una secuencia dada mediante redes de flujo.

\subsection{Modelo de Programaci\'on Entera Mixta (Beasley et al., 2000)\cite{beasley2000}}

La formulaci\'on de Beasley et al. \cite{beasley2000} es el est\'andar para el ALP est\'atico en entornos de m\'ultiples pistas. El modelo integra decisiones l\'ogicas (asignaci\'on y orden) con variables continuas (tiempo).

\begin{table}[htbp]
\centering
\renewcommand{\arraystretch}{1.2}
\begin{tabular}{|l|p{11.5cm}|}
\hline
\textbf{Notaci\'on} & \textbf{Par\'ametro} \\
\hline
$A$ & Conjunto de aeronaves. \\
$K$ & Conjunto de pistas. \\
$T_i$ & Tiempo de aterrizaje objetivo o preferido para la aeronave $i$. \\
$[E_i, L_i]$ & Ventana de tiempo permitida para el aterrizaje ($L_i > E_i$). \\
$c_i^-$ & Costo de penalizaci\'on ($\geq 0$) por unidad de tiempo por aterrizar antes de $T_i$. \\
$c_i^+$ & Costo de penalizaci\'on ($\geq 0$) por unidad de tiempo por aterrizar despu\'es de $T_i$. \\
$S_{ij}$ & Separaci\'on por turbulencia de estela entre $i$ y $j$, si aterrizan en la \textbf{misma pista} ($i$ antes que $j$). \\
$s_{ij}$ & Separaci\'on requerida entre $i$ y $j$, si aterrizan en \textbf{pistas distintas} ($i$ antes que $j$). \\
\hline
\end{tabular}
\caption{Par\'ametros de entrada del modelo de Beasley \cite{beasley2000}.}
\end{table}

\noindent \textbf{Variables de Decisi\'on}\\
El modelo utiliza variables binarias para la asignaci\'on y secuencia, y variables continuas para los tiempos:

\begin{itemize}
    \item $a_{ik} = 1$ si la aeronave $i$ es asignada a la pista $k$; $0$ en caso contrario.
    \item $z_{ij} = 1$ si las aeronaves $i$ y $j$ se asignan a la \textbf{misma pista}; $0$ en caso contrario.
    \item $\delta_{ij} = 1$ si la aeronave $i$ aterriza antes que la $j$; $0$ en caso contrario.
    \item $x_i$: Tiempo real de aterrizaje asignado a la aeronave $i$.
    \item $x_i^-, x_i^+$: Variables de desviaci\'on para linealizar la funci\'on objetivo, donde $x_i^- = \max(0, T_i - x_i)$ y $x_i^+ = \max(0, x_i - T_i)$.
\end{itemize}

\noindent \textbf{Formulaci\'on Matem\'atica (MIP)}
\begin{alignat}{2}
& \underset{a,z,\delta,x,x_i^-,x_i^+}{\textbf{Minimizar}} \qquad & \sum_{i \in \mathcal{A}} (c_i^- x_i^- + c_i^+ x_i^+) & \label{eq:beasley_obj} \\
& \textbf{Sujeto a:} & \nonumber \\
& & E_i \leq x_i \leq L_i & \qquad \forall i \in \mathcal{A} \\
& & 0 \leq x_i^- \leq T_i - E_i & \qquad \forall i \in \mathcal{A} \\
& & 0 \leq x_i^+ \leq L_i - T_i & \qquad \forall i \in \mathcal{A} \\
& & x_i = T_i - x_i^- + x_i^+ & \qquad \forall i \in \mathcal{A} \\
& & \delta_{ij} + \delta_{ji} = 1 & \qquad \forall i, j \in \mathcal{A} : i \neq j \\
& & \sum_{k \in \mathcal{K}} a_{ik} = 1 & \qquad \forall i \in \mathcal{A} \\
& & z_{ij} = z_{ji} & \qquad \forall i, j \in \mathcal{A} : i < j \\
& & z_{ij} \geq a_{ik} + a_{jk} - 1 & \qquad \forall i, j \in \mathcal{A} : i < j, k \in \mathcal{K} \\
& & x_j \geq x_i + S_{ij}z_{ij} + s_{ij}(1 - z_{ij}) - M\delta_{ji} & \qquad \forall i, j \in \mathcal{A} : i \neq j \label{eq:beasley_sep} \\
& & a_{ik}, z_{ij}, \delta_{ij} \in \{0, 1\} & \qquad \forall i, j, k
\end{alignat}

\noindent \textit{Nota sobre el caso de Pista \'Unica:} Para los prop\'ositos de este proyecto, donde se considera una sola pista ($\mathcal{K} = \{1\}$), el modelo se simplifica dr\'asticamente. Las variables $a_{i1}$ y $z_{ij}$ siempre toman el valor de $1$, lo que elimina la necesidad de las restricciones de asignaci\'on y reduce la restricci\'on de separaci\'on \cite{beasley2000} a la forma cl\'asica de pista \'unica: $x_j \geq x_i + S_{ij} - M\delta_{ji}$.

\section{Representaci\'on y Manejo de Restricciones}
Para abordar el problema ARSP la representación escogida fue un arreglo de permutación. Esto debido a su uso estandarizado en la literatura del problema, esta representación modela la secuencialidad inherente al problema.

\subsection{Manejo de Restricciones}
El ARSP tiene restricciones duras que invalidan una amplia porción del espacio de Búsqueda. El algoritmo propuesto maneja estas limitantes en la Función de Evaluación, de la siguiente forma:
\begin{itemize}
    \item \textbf{Ventanas de Tiempo Operativas($L_i$):} Mientras se evalúa una permutación, en caso de que el tiempo proyectado para el aterrizaje de una aeronave exceda el tiempo máximo permitido de vuelo, si hablamos en un escenario de la vida real esto implica que el avión se quedará sin combustible antes de poder aterrizar, la solución candidata se etiqueta como infactible (retornando un valor arbitrario de -1), descartando ese vecino
    \item \textbf{Estela de Turbulencia Asimétrica($S_{ij}$):} Para que el aterrizaje sea seguro, el algoritmo compara dinámicamente a la aeronave actual contra todos los aviones que hayan aterrizado previamente, osea que estén en posiciones anteriores en la permutación que se está revisando. Considerando los tiempos de aterrizaje de estos y la matriz de separación, se calcula  el tiempo mínimo seguro en el que la pista estará liberada.
\end{itemize}

\section{Descripci\'on del algoritmo}
El algoritmo metaheurístico implementado corresponde a una búsqueda local basada en trayectorias: \textbf{Hill Climbing con Mejor Mejora y Restarts (HC-BI)}. Esta técnica se clasifica como un algoritmo de búsqueda incompleta, de tipo iterativo, que evalúa el entorno local que rodea a una solución candidata, buscando minimizar el valor de la función de evaluación, que corresponde al coste de la permutación

\subsection{Operador de Movimiento (\textit{Swap})}
Para poder movernos en el espacio de búsqueda se determinó el uso del movimiento swap. Este operador selecciona dos posiciones distintas dentro de la permutación e intercambia las aeronaves de esas posiciones, con lo que no se pierden soluciones y no se repiten naves, por lo que la representación no se rompe. Como estamos en la versión Best Improvement de Hill Climbing, el algoritmo para esto es exhaustivo, cada vez que se genere el vecindario se generan todas las combinaciones posibles de pares de posiciones.

A modo de ejemplo gráfico, partiendo de la permutación inicial $[1, 2, 3]$, el vecindario completo generado por el movimiento \textit{Swap} sería:
\begin{itemize}
    \item Intercambio posiciones 0 y 1 $\rightarrow [2, 1, 3]$
    \item Intercambio posiciones 0 y 2 $\rightarrow [3, 2, 1]$
    \item Intercambio posiciones 1 y 2 $\rightarrow [1, 3, 2]$
\end{itemize}

\subsection{Procedimiento General HC-BI}
El algoritmo pide una solución factible como solución inicial para poder tener un desempeño que sea considerado útil, ya que puede estancarse fácilmente en otro escenario. A continuación, se presenta el pseudocódigo de su funcionamiento:

\begin{verbatim}
Procedimiento HC_BI_Restarts():
    Mientras num_prueba < max_restarts Hacer:
        local = Falso
        s_c = Generar permutación estocástica (aleatoria)
        
        // Búsqueda de una Solución Inicial Factible
        Mientras f(s_c) == Inviable Hacer:
            Cambiar semilla estocástica
            s_c = Generar nueva permutación estocástica
        
        // Ciclo Hill Climbing
        Mientras local == Falso Hacer:
            Generar vecindario completo N(s_c) mediante movimiento Swap
            s_n = Seleccionar el mejor vecino factible dentro de N(s_c)
            
            Si f(s_n) < f(s_c) Entonces:
                s_c = s_n
            Sino:
                local = Verdadero
        
        Actualizar Mejor_Solucion_Global si f(s_c) es el menor histórico
        num_prueba = num_prueba + 1
\end{verbatim}

\subsection{Análisis de la ejecución paso a paso:}
\begin{enumerate}
    \item \textbf{Control de Restarts:} El bucle externo controla la repetición del algoritmo hasta alcanzar el límite máximo de iteraciones, de restarts, permitidos.
    \item \textbf{Inicialización Factible Estocástica:} Siguiendo las indicaciones para la técnica, se genera una solución inicial aleatoria. Ahora, para que la solución no inicie en un, más que posible, valle de infactibilidad, lo que se hace es cambiar la semilla para generar una nueva permutación que sea factible.
    \item \textbf{Exploración Exhaustiva (Best Improvement):} A partir de la solución actual, se genera todo el vecindario mediante el operador swap. Se hace la evaluación matemática de cada vecino y se registra la Mejor Mejora según la Función de Evaluación, con eso se hace el salto a la nueva solución actual.
    \item \textbf{Convergencia y Reinicio:} El ciclo de mejoras continúa sin parar, hasta que ningún vecino logre ser mejor que la solución actual. Ahí se determina que el algoritmo está en un óptimo local,se actualiza la mejor solución global, si se da el caso, y se genera un reinicio.
\end{enumerate}

\section{Experimentos}
El objetivo de la experimentación es evaluar tanto el rendimiento como la capacidad de convergencia del Algoritmo HC-BI en dos situaciones del problema: Una Pista y Varias pistas (de 2 a 5)

\subsection{Entorno de Ejecución}
Para garantizar la eficiencia computacional al evaluar vecindarios masivos, el algoritmo fue implementado en C++ (estándar \texttt{C++11}) y compilado utilizando \texttt{g++}. Las simulaciones se ejecutaron en el subsistema Windows Subsystem for Linux (WSL), sobre un equipo equipado con un procesador AMD Ryzen 5 8500G y 32 GB de memoria RAM.

\subsection{Instancias de Prueba y Criterios de Exclusión}
El algoritmo fue ejecutado sobre dos conjuntos de datos:
\begin{itemize}
    \item \textbf{Instancias Unipista (Beasley):} Se seleccionaron 12 de las 13 instancias clásicas de la OR-Library, de \texttt{airland1} a \texttt{airland12}. La instancia 13 fue excluida luego de iniciar los experimentos debido a una explosión combinatoria que elevó el tiempo de ejecución a más de 30 minutos para una sola semilla, este tiempo masivo vuelve inviable la realización de pruebas repetidas en un marco de tiempo razonable.
    \item \textbf{Instancias Multipista (FTP):} Se evaluaron las 12 instancias disponibles, para cada una de estas se le hicieron pruebas con distintas cantidades de pistas, en busca de simular la variabilidad de un caso real, evaluando escenarios con $K\in\{2,3,4,5\}$ pistas.
\end{itemize}

\subsection{Parámetros y Variables de Estudio}
Para analizar el comportamiento del algoritmo ante distintos niveles de diversificación se definieron dos parámetros que cambiamos para contrastar la posible mejora que obtenemos de diversificar más.
\begin{itemize}
    \item \textbf{Limite de Restarts (10 vs 20):} Cuando el algoritmo alcanza un óptimo local el algoritmo llama a un reinicio, en este punto hace un cambio en la seed para conseguir una solución inicial distinta. Se experimentó limitando la cantidad de Restarts a 10 y luego a 20. El punto de este aumento es observar el \textit{trade-off} (compromiso) que existe entre el tiempo de ejecución y la calidad de la solución final, buscamos determinar si duplicar la exploración del espacio justifica el costo temporal copartícipe.
    \item \textbf{Robustez Estocástica (5 vs 10):} Para evitar posibles sesgos ligados a una inicialización fortuita, se probaron grupos de 5 y 10 semillas iniciales distintas. Esto permite verificar si la convergencia del algoritmo es estadísticamente robusta (baja varianza en los resultados) independiente de la solución inicial dada.
\end{itemize}

\subsection{Métricas de Evaluación}
Para comparar el posible impacto de los parámetros mencionados, los resultados que se presentan en la siguiente sección son analizados en base a 3 métricas clave:
\begin{itemize}
    \item \textbf{Mejor Costo Global (Best):} El valor mínimo de penalización alcanzado en todas las ejecuciones
    \item \textbf{Costo Promedio (Avg) y Varianza:} Para establecer si el método heurístico es estable
    \item \textbf{Tiempo de Ejecución (ms):} Esta métrica es para determinar la aplicabilidad del algoritmo en entornos realistas, osea, si es escalable.
\end{itemize}

\section{Resultados}
En esta sección se presentan los resultados computacionales obtenidos tras someter el algoritmo HC-BI a las distintas instancias de prueba. Para medir el impacto real de la heurística, los resultados se comparan contra una política de línea base \textit{First-Come, First-Served} (FCFS), la cual representa el orden natural de llegada de las aeronaves sin ningún tipo de optimización inteligente.

\subsection{Resultados en Pista Única (Instancias Airland)}
En la Tabla \ref{tab:unipista} se resume el desempeño del algoritmo para las instancias \texttt{airland1} a \texttt{airland12}. Se utilizó la configuración más robusta (ej. 10 semillas y 20 restarts) para extraer el mejor costo posible.

% TABLA 1: UNIPISTA VS FCFS
\begin{table}[htbp]
\centering
\caption{Comparación de Costos: HC-BI vs Política FCFS en Pista Única}
\label{tab:unipista}
\begin{tabular}{|l|c|c|c|c|c|}
\hline
\textbf{Instancia} & \textbf{Costo FCFS} & \textbf{Mejor Costo HC} & \textbf{Costo Prom} & \textbf{Mejora (GAP \%)} & \textbf{Tiempo Prom (ms)} \\ \hline
airland1 & 1233.0 & 1150.0 & 1150.0 & 6.7\% & 0.3 \\ \hline
airland2 & 2063.0 & 1720.0 & 1875.0 & 16.6\% & 0.7 \\ \hline
airland3 & 2940.0 & 1610.0 & 1610.0 & 45.2\% & 4.7 \\ \hline
airland4 & 4566.0 & 4480.0 & 4480.0 & 1.8\% & 3.2 \\ \hline
airland5 & 7232.0 & 4800.0 & 4800.0 & 33.6\% & 3.2 \\ \hline
airland6 & 24442.0 & 24442.0 & 24442.0 & 0.0\% & 1.4 \\ \hline
airland7 & 4004.0 & 3974.0 & 3974.0 & 0.8\% & 3.1 \\ \hline
airland8 & 4669.0 & 3240.0 & 3553.0 & 30.6\% & 3502.5 \\ \hline
airland9 & 17552.9 & 10345.0 & 10468.0 & 41.1\% & 3782.8 \\ \hline
airland10 & 41623.1 & 25843.4 & 26561.3 & 37.9\% & 16540.2 \\ \hline
airland11 & 33888.4 & 22522.5 & 22770.4 & 33.5\% & 57810.4 \\ \hline
airland12 & 99409.0 & 27943.6 & 28009.2 & 71.9\% & 177177.1 \\ \hline
\hline
\textbf{Promedio} & \textbf{20301.9} & \textbf{11141.1} & \textbf{11137.7} & \textbf{26.7\%} & \textbf{21569.1} \\ \hline
\end{tabular}
\end{table}

\newpage
\textbf{Análisis Unipista:} Como se observa en la Tabla \ref{tab:unipista}, el algoritmo HC-BI logra superar ampliamente a la regla FCFS en todas las instancias. Las mejoras porcentuales oscilan entre un [MIN GAP]\% y un espectacular [MAX GAP]\%. Resulta interesante notar el caso de \texttt{airland6}, donde la alta restricción de la ventana de tiempo ($L_i$) obligó a usar una inicialización estocástica asistida, logrando tiempos de ejecución sumamente competitivos ([TIEMPO AIRLAND6] ms) dada la complejidad del vecindario evaluado.


\subsection{Escalabilidad en Múltiples Pistas (Instancias FPT)}
El segundo experimento evaluó la capacidad de la asignación \textit{Greedy} en base a una permutación estocástica para distribuir la carga de tráfico al tener habilitadas 2,3,4 y 5 pistas. La Tabla \ref{tab:multipista} muestra las penalizaciones conforme crece la infraestructura.

% TABLA 2: MULTIPISTA (ESCALABILIDAD)
\begin{table}[h]
\centering
\caption{Evolución del Mejor Costo (HC-BI) al aumentar la cantidad de pistas ($K$)}
\label{tab:multipista}
\begin{tabular}{|l|c|c|c|c|}
\hline
\textbf{Instancia} & \textbf{K = 2 Pistas} & \textbf{K = 3 Pistas} & \textbf{K = 4 Pistas} & \textbf{K = 5 Pistas} \\ \hline
FPT01 & 177222 & 107607 & 73127 & 52932 \\ \hline
FPT02 & 186436 & 112885 & 76359 & 54844 \\ \hline
FPT03 & 190342 & 113378 & 75248 & 54044 \\ \hline
FPT04 & 194645 & 117921 & 79757 & 57208 \\ \hline
FPT05 & 194365 & 115738 & 76425 & 53092 \\ \hline
FPT06 & 178397 & 107179 & 72499 & 51727 \\ \hline
FPT07 & 182076 & 108829 & 72379 & 50654 \\ \hline
FPT08 & 160641 & 96179  & 64698 & 45960 \\ \hline
FPT09 & 179699 & 107410 & 71707 & 50404 \\ \hline
FPT10 & 189005 & 114430 & 77883 & 55987 \\ \hline
FPT11 & 194438 & 115862 & 76855 & 53728 \\ \hline
FPT12 & 211461 & 123736 & 80959 & 55771 \\ \hline
\hline
\end{tabular}
\end{table}

\textbf{Análisis Multipista:} Los resultados validan empíricamente el impacto de la infraestructura en la puntualidad. Al pasar de 2 a 3 pistas, se observa una reducción promedio de 40.1\% en la penalidad. Si se siguieran aumentando las pistas, llegará el momento donde los costes pueden llegar a 0, lo que indicaría que el aeropuerto estaría sobredimensionado para esa carga de tráfico y todas las aeronaves logran aterrizar en su tiempo objetivo ($T_i$) sin generar multas, un caso ideal, pero inviable económicamente.


\subsection{Análisis Paramétrico: El Trade-off de Exploración vs Explotación}
Finalmente, se evaluó cómo influye el aumento del esfuerzo computacional (pasar de 10 a 20 \textit{restarts} y de 5 a 10 semillas) en la calidad de la solución.

\textbf{Análisis de Restarts:} Al analizar los datos consolidados, se determinó que duplicar la cantidad de \textit{restarts} de 10 a 20 generó un incremento lineal en el tiempo de cómputo (aumentando un 14.3\% en promedio). Sin embargo, la mejora en el \textit{Mejor Costo Global} fue nula. Esto demuestra que la heurística HC-BI implementada posee una alta tasa de convergencia temprana: encuentra valles de muy alta calidad en sus primeras iteraciones y agotar más restarts solo incurre en gasto de CPU sin aportar beneficio alguno al control de tráfico aéreo.

\section{Conclusiones}
%Conclusiones RELEVANTES del estudio realizado. Incluir conclusiones acerca de la adecuaci\'on de la propuesta de soluci\'on al problema que se busca resolver. Listar y analizar ventajas y desventajas de la propuesta en base a los resultados obtenidos y comportamiento de la propuesta en diferentes escenarios (problemas/instancias/par\'ametros). Incluir trabajo futuro en base a las conclusiones.
El presente informe abordó el Problema ARSP, un problema central para la eficiencia de la industria aeorportuaria y para la reducción de costos operativos. Se implementó una maetaheurística de búsqueda local (HC-BI), programada en C++, que permitió superar las barreras de tiempos inviables que erigen los métodos exactos.

La experimentación en escenarios de pista única demostró la superioridad de la optimización mediante heurísticas frente a la metodología FCFS, o FIFO también llamada. El algoritmo HC-BI logró una reducción de costos promedio del 26.5\% sobre FCFS, alcanzando mejoras máximas de hasta un 71.9\% en instancias de alta complejidad como \texttt{airland12}.
Además, el algoritmo implementado manejó las restricciones duras de seguridad, descartando soluciones infactibles que excedían el tiempo máximo ($L_i$) y respetando las separaciones necesarias por la estela de turbulencia($S_{ij}$).

En cuanto a la escalabilidad, las pruebas con instancias FTP en múltiples pistas demostraron que la asignación \textit{Greedy} logra distribuir eficazmente la carga de tráfico. Al escalar de 2 a 3 pistas disponibles se apreció una mejora promedio del 40.1\% en la penalización. Estos resultados demuestran el gran impacto de la infraestructura en la puntualidad, lo que nos permite buscar matemáticamente el momento en que el aeropuerto se vuelve sobredimensionado y las penalizaciones desaparecen.

Por último, el estudio con respecto a la variación de parámetros reveló una característica del algoritmo, tiene una convergencia temprana alta. Cuando aumentamos los \textit{Restarts} de 10 a 20 el tiempo de procesamiento aumentó un 14.3\%, mas la mejora que se buscaba resultó ser nula. Esto prueba que el forzar más reinicios, buscando la diversificación, solo generó un uso de CPU que no se justificó debido a los nulos beneficios que reportó. Por ende, se puede concluir que el diseño heurístico propuesto, en su configuración base (5 semillas, 10 restarts) se presenta como una herramienta rápida, robusta y computacionalmente eficiente, viable para entornos operacionales que exigen respuestas de calidad en tiempos muy pequeños.

\section{Uso de LLMs e Inteligencia Artificial}
A continuaci\'on, se detallan las herramientas de Inteligencia Artificial y Modelos de Lenguaje Grande (LLMs) utilizadas como apoyo en el desarrollo y redacci\'on del presente informe, en cumplimiento con las normativas de probidad acad\'emica de la asignatura.

\vspace{0.3cm}
\noindent \textbf{Herramienta 1: Asistencia en Redacci\'on y Formato}
\begin{enumerate}
    \item \textbf{Herramienta utilizada:} Gemini (Google) - Versi\'on Advanced (Modelo 3.1 Pro Extendido).
    \item \textbf{Prop\'osito de uso:} Estructuraci\'on general del informe, asistencia en la traducci\'on de conceptos t\'ecnicos, depuraci\'on del formato y c\'odigo de ecuaciones en \LaTeX{}, y s\'intesis narrativa de la revisi\'on bibliogr\'afica.
    \item \textbf{Prompts relevantes:} Debido a la metodolog\'ia de conversaci\'on continua y guiada, se omite la transcripci\'on de \textit{prompts} individuales. En su lugar, el registro \'integro de las interacciones, directrices y la evoluci\'on del documento puede ser auditado directamente en el siguiente enlace al historial del chat: \textbf{[https://gemini.google.com/share/37bac553384f]}.
    \item \textbf{Nivel de edici\'on:} Altamente iterativo. La herramienta se utiliz\'o para generar primeros borradores (\textit{drafts}) que luego fueron analizados, parafraseados y reescritos manualmente para asegurar el hilo conductor y el tono acad\'emico adecuado. Ninguna secci\'on fue copiada y pegada ciegamente; el c\'odigo \LaTeX{} devuelto fue depurado manualmente.
    \item \textbf{Validaci\'on:} Toda afirmaci\'on matem\'atica o te\'orica generada fue contrastada mediante la lectura cruzada con los documentos en PDF originales (ej. Beasley 2000, Cao \& Xu 2024), adem\'as de verificar el estricto cumplimiento con la r\'ubrica de evaluaci\'on.
\end{enumerate}

\noindent \textbf{Herramienta 2: Automatización de Datos y Scripts de Procesamiento}
\begin{itemize}
    \item \textbf{Herramienta utilizada:} Gemini (Google) - Versi\'on Advanced (Modelo 3.1 Pro Extendido).
    \item \textbf{Propósito:} Generación de un script de \texttt{Python} (\texttt{extractor.py}) utilizando expresiones regulares para parsear masivamente más de 120 archivos \texttt{.txt} anidados en el subsistema Linux (WSL). El script extrajo parámetros clave (costos y tiempos en milisegundos) consolidándolos en un CSV unificado. Adicionalmente, se utilizó para la formulación de funciones lógicas y matriciales en Google Sheets para el cálculo automático de promedios, desviaciones estándar y el GAP porcentual frente al baseline FCFS.
    \item \textbf{Prompts Relevantes:} \textit{"Crea un script de python para tomar todos los datos de tiempo, penalización, los datos que nos importan al fin y al cabo para llenar el excel"} y \textit{"Dame la función de Excel para que se hagan los calculos de manera automática, considera que estoy en google spreadsheets"}.
    \item \textbf{Nivel de Edición:} El código base proporcionado por la IA fue adaptado localmente para integrarse de forma dinámica con la estructura de rutas relativas y subcarpetas específicas del entorno de desarrollo del proyecto.
    \item \textbf{Validación:} El archivo CSV generado y los cálculos de las hojas de cálculo fueron verificados manualmente, realizando un muestreo aleatorio cruzado entre los archivos de texto originales arrojados por el código fuente en C++ y las celdas finales del resumen, garantizando un 100\% de fidelidad en la tabulación de datos.
\end{itemize}

%------Esta sección puede ser removida ----------%
\section{Bibliograf\'ia}
Indicando toda la informaci\'on necesaria de acuerdo al tipo de documento revisado. Todas las referencias deben ser citadas en el documento.
%------Remover hasta acá ----------%

\bibliographystyle{plain}
\bibliography{Referencias}

\end{document} 
